\section{Conclusion}\label{sec:conclusion}

We have described a text-associated multi-morphology motion corpus and a
physics-based imitation system that preserves the correspondence between
source clips, generated references, and simulated bodies. The construction
combines AMASS/HumanML3D provenance, HUMOS generation, SMPLSim assets,
conservative reference refinement, and deterministic clip-grouped splits.
The controller uses temporal slot/type attention with morphology information
in both actor and critic and is trained directly across body configurations.

Offline evaluation of the selected checkpoint gives 98.47\% success on
validation and 99.14\% on the untouched test split for the one-shape-per-clip
trained-body protocol, and the eight-run development grid selects slot/type
attention without adaLN-Zero. Residual failures on both splits are dominated
by ground-contact transitions and whole-body acrobatics. Replaying this
checkpoint across all 19,200 (clip, shape) pairs of the development corpus
finds no statistically detectable evidence that dynamic motions are more
shape-sensitive to body morphology than quasi-static ones, once a
near-zero-baseline normalization artifact in the naive comparison is
corrected -- a properly negative result rather than an unresolved question.
These results do not establish all-shape coverage or unseen-body
generalization: an evaluation on body shapes inside the trained coefficient
range but excluded from the 128 training configurations remains future work.
